\section{A coefficient Hom countertest and its vanishing standard-target residue}
\label{sec:136-q3-coefficient-test}

\paragraph{Scope.}
The archive locators are cote136, physical pp.~117,119,121,
footers~116,118,120, scans~117,119,121.  The open issue is the
degree-three graph-linking condition after a free degree-three
Tate stage.  We use the complete checked arity filtration
ARCH-136-Q3-LOW-FILTRATION.  The first assertion below excludes
one proposed module-theoretic shortcut.  The same explicit
algebra does not give an actual graph-residue counterexample
under the standard-target hypotheses stated in the second
assertion.  No general preservation theorem is deduced.

\paragraph{ARCH-136-Q3-COEFFICIENT-TEST.}
Put $S=\Gamma_{\mathbb Z}(T)$, $J=\mathbb Q[T]/S$, and
$\bar S=S/2$.  Let $h=2^a$, $x=T^{[h]}$, and write
$\bar x$ for its reduction.  Set
$I=\operatorname{Ann}_{\bar S}(\bar x)$.
Then
\[
 I=(\bar x),\qquad I^2=0,
 \qquad
 \operatorname{Hom}_{S,\mathrm{gr}}(I\langle6h\rangle,J)\ne0.
 \tag{136.CT1}
\]
There is such a homomorphism nonzero in arbitrarily high
output weights.  In particular neither graded injectivity
of $J$ nor multiplication by fresh coefficient weights proves
vanishing for all maps from $(I\cap I)/I^2$ to $J$.

This exact coefficient kernel occurs in the actual arity
filtration of the following free algebra.  Start with
\[
 A=S\langle v_1,v_2,u_1,u_2\rangle,\quad
 |v_i|=2,\ |u_i|=3,\quad
 \mathrm{wt}(v_i)=\mathrm{wt}(u_i)=2h,\quad dv_i=2u_i.
 \tag{136.CT2}
\]
Attach $|c_i|=2$, $|t_i|=1$, of weight $3h$, with
\[
 dc_i=xu_i,\qquad dt_i=xv_i-2c_i.
 \tag{136.CT3}
\]
Write $y=T^{[2h]}$ and $x^2=2\kappa y$, where
$\kappa=\binom{2h}{h}/2$ is odd.  The three-form
\[
 \Omega=x(t_1c_2-c_1t_2)
                  +\kappa y(v_1t_2-v_2t_1)
 \tag{136.CT4}
\]
is closed and represents a nonzero class of exact order two,
with order relation
\[
                         2\Omega=d(x t_1t_2).
 \tag{136.CT5}
\]
Its arity-two component generates the coefficient ideal $I$
in the surviving associated quotient.

Now work in the standard integral theory $B$ on
$(X,G)=(\Delta[2],\partial\Delta[2])$.  Choose $h$ sufficiently
large for its relative coefficient comparison through form
degree two in weights $2h$ and $3h$.  Require all displayed
generators and pairs to have zero restriction to $G$.
Let the reference map send all old generators in (136.CT2)
to zero.  Let a second map send $u_i$ to zero and $v_i$ to
relative two-forms $\zeta_i$ whose integral periods are even.
The latter is exactly the old graph-linking condition for
$[u_i]$ against this reference map.  Then the graph residue
of $[\Omega]$, and of every coefficient multiple of it, is
zero for every extension of these maps through (136.CT3)
with the prescribed zero boundary.  Thus (136.CT1) does not
by itself produce a nonzero realized graph residue.

\paragraph{Proof of the module and algebra assertions.}
The binary-digit basis gives
\[
       \bar S\simeq
          \bigotimes_{r\ge0}\mathbb F_2[x_r]/(x_r^2),
                    \qquad x_r=T^{[2^r]}\bmod2.
 \tag{136.CT6}
\]
Indeed a product of distinct digit divided powers has odd
multinomial coefficient and gives the divided-power basis
element with that support, whereas a repeated digit gives an
even coefficient.  These parity statements follow directly
from $\nu_2(n!)=n-s_2(n)$, obtained by summing the numbers
of multiples of $2,4,8,\ldots$; here $s_2(n)$ is the number
of binary digits equal to one.  Thus (136.CT6) is a basis
calculation, not a finiteness assumption.
Split off the $a$-th factor, writing
$\bar S=\bar S'\oplus\bar x\bar S'$.
Multiplication by $\bar x$ carries the first summand
isomorphically to the second and kills the second.  This
proves $I=(\bar x)$ and $I^2=0$.

For any $v\ge0$ with zero $a$-th binary digit, define
\[
 I\langle v\rangle\longrightarrow J,
                  \qquad r\longmapsto rT^{[v]}/2\pmod S.
 \tag{136.CT7}
\]
Integral lifts of $r$ differing by twice an element of $S$
give the same result; hence this is a graded $S$-linear map.
Its value on $\bar x$ is nonzero, since
$\binom{h+v}{h}$ is odd by the same digit calculation.
Choose $v=6h$, which has zero $a$-th digit.  Multiplying
$\bar x$ by a digit $T^{[2^b]}$ beyond the supports of $h$
and $6h$ still gives a nonzero value in (136.CT7).  These
weights are unbounded.  This proves (136.CT1).

In (136.CT2), listing monomials in forms at most three gives
\[
 H^0(A)=S,\quad H^1(A)=H^2(A)=0,\quad
 H^3(A)=(S/2)[u_1]\oplus(S/2)[u_2].
 \tag{136.CT8}
\]
The two old third classes are independent.  In the first
differential of the checked arity filtration, a coefficient
$s$ of $\tau=t_1c_2-c_1t_2$ survives precisely when both
$s\bar x[u_i]$ vanish.  Its coefficient kernel is therefore
exactly $I$.
For the actual closed lift, compute
\[
 \begin{aligned}
 d\tau&=x(u_2t_1-u_1t_2)+x(v_1c_2-v_2c_1),\\
 d(v_1t_2-v_2t_1)
       &=2(u_1t_2-u_2t_1)-2(v_1c_2-v_2c_1),\\
 d(t_1t_2)&=2\tau+x(v_1t_2-v_2t_1).
 \end{aligned}
 \tag{136.CT9}
\]
The two even terms $xv_1v_2$ cancel.  Substituting
$x^2=2\kappa y$ proves (136.CT4)--(136.CT5).
The nonzero top component $\bar x\tau$ survives, by this
actual lift.  There is no incoming differential to it by
ARCH-136-Q3-LOW-FILTRATION, so $[\Omega]\ne0$ and its
order is exactly two.  Its coefficient multiples cover
$I$ in the associated quotient.  This does not claim a
splitting of the whole cohomology filtration.

\paragraph{Proof of the standard-target assertion.}
In the stated weights the standard relative comparison gives
$H^1(B(X,G))=0$ and $H^2(B(X,G))=S^w$, with the canonical
period.  Its relative form groups are integer-torsion-free,
and $B_2$ has no forms of degree three.  The long exact
sequence of reduction modulo two therefore gives
\[
 H^1(B(X,G)/2)=0,\qquad H^2(B(X,G)/2)=S^w/2.
 \tag{136.CT10}
\]
The even-period hypothesis makes $\zeta_i$ exact in the
second complex.  Lift a primitive to an integral relative
one-form $r_i$, and write
\[
                   \zeta_i=dr_i+2\theta_i
 \tag{136.CT11}
\]
with an integral relative two-form $\theta_i$.
The assignment
\[
                       t_i\longmapsto xr_i,
                \qquad c_i\longmapsto x\theta_i
 \tag{136.CT12}
\]
is a strict extension: $dc_i=0$ in the target, and its other
differential is $x\zeta_i-2x\theta_i$ as required.
It preserves all the zero boundary data.  The image of
$\Omega$ is zero by form degree, and its order relation has
image
\[
               x^3r_1r_2=2x\kappa y r_1r_2.
 \tag{136.CT13}
\]
The period of this relation divided by two is integral.
This is exactly zero graph residue against the zero
reference extension.

It remains to check that different allowed extensions do
not change this conclusion.  Their one-form differences
$\delta t_i$ are relative and have differential divisible
by two.  The first vanishing in (136.CT10) gives
\[
                   \delta t_i=df_i+2h_i
 \tag{136.CT14}
\]
with relative zero-forms $f_i$ and relative one-forms $h_i$.
Modulo two, the change in the period of $x t_1t_2$ is
the sum of the periods of
\[
        x\,df_1t_2,\qquad x\,t_1df_2,
                               \qquad x\,df_1df_2.
 \tag{136.CT15}
\]
Stokes has no boundary term because the $f_i$ are relative.
The first two periods are, with their respective signs,
periods of $x f_i dt_j$.  Modulo two one has
$dt_j=x\zeta_j$, so they vanish because $x^2=0$ modulo
two.  The third term is exact relative to $G$ and its
period is zero.  Thus the order-relation period remains
even for every extension.  The same applies to the
reference map and to all coefficient multiples, by
$S$-linearity of the period.  This proves the statement.
\hfill$\square$

\paragraph{Exact remaining issue.}
The module Hom is nonzero, but this concrete first test has
zero actual residue under its explicit standard-target and
zero-reference hypotheses.  These calculations do not control the general mixed
$\ker\Phi$ piece or finite relation matrices in $\ker d_2$.  In particular, an arbitrary admissible
target, arbitrary old maps, or a whole repair iteration is
not covered by the second assertion.

\paragraph{Verification status.}
The author completed a whole saved-proof reading: the binary
coefficient presentation and graded Hom, unbounded output weights,
complete old low cohomology, every differential in (136.CT9),
the exact order-two class, relative reduction modulo two, all
target choices and Stokes signs were checked.  No gap was found.
The root also read the complete saved proof and passed it at
the stated standard-target scope. The separate events are in
the register. No main-gate status change is claimed.
